% ============================================================================
% THE FIELD: EMERGENT ATTENTION THROUGH EIGENMODE SELECTION
% ============================================================================
\subsection{The Field: Emergent Attention via Eigenmode Selection}
\label{sec:field}

The Field is not a data structure that nodes query.  It is a
\emph{force that acts on them}---analogous to a physical field acting on
particles.  It emerges from the gossip protocol
(\Cref{sec:gossip}): as nodes exchange \texttt{FieldDigest} messages,
each node's \texttt{FieldView} detects emergent clusters of
structurally aligned, phase-coherent tries, identifies the dominant
cluster, and projects asymmetric pressure back onto the local trie.
This mechanism functions as the system's attention: it selects what to
retain and what to forget, without matrices, softmax, or learned
parameters.

\subsubsection{Affine Coupling: Multi-Scale Structural Alignment}

Structural similarity between two nodes is measured by
\emph{affine coupling}, which evaluates overlap between their
512-bit affinity vectors at multiple spatial scales.  For affinity
vectors $\mathbf{a}, \mathbf{b} \in \{0,1\}^{8 \times 64}$ (eight
64-bit words), the coupling is:
\begin{equation}
  C_{\text{aff}}(\mathbf{a}, \mathbf{b})
    = \frac{1}{8 \cdot 64 \cdot |\mathcal{K}|}
      \sum_{k \in \mathcal{K}} \sum_{w=0}^{7}
      \text{popcount}\!\bigl(\pi_k(\mathbf{a})_w \;\mathbin{\&}\; \mathbf{b}_w\bigr)
  \label{eq:affine_coupling}
\end{equation}
where $\mathcal{K} = \{1, 4, 16, 64\}$ are the rotation exponents
(base-4 exponential spacing: $4^0, 4^1, 4^2, 4^3$), and $\pi_k$ is an
affine permutation that rotates word indices by $k$ and bit positions
within each word by $(3k) \bmod 64$:
\begin{equation}
  \pi_k(\mathbf{a})_w
    = \texttt{rotl}\!\bigl(\mathbf{a}_{(w+k) \bmod 8},\; (3k) \bmod 64\bigr)
  \label{eq:affine_permutation}
\end{equation}

The multiplier 3 is chosen because $\gcd(3, 64) = 1$, ensuring that
the composed shift visits many distinct bit phases as $k$ sweeps the
exponential scale.  The resulting shifts $\{3, 12, 48, 0\} \bmod 64$
span the full 64-bit lane.

Affine coupling returns a value in $[0, 1]$ where $1$ indicates
perfect alignment at all scales.  The coupling threshold for eigenmode
membership is $\theta_C = 0.15$, deliberately below the chance-overlap
level for random 512-bit vectors ($\approx 0.25$) because eigenmode
assignment is a two-stage gate: nodes must also pass the phase
coherence check.

\subsubsection{Phase Coherence: Surprisal Velocity Alignment}

Structural alignment alone is insufficient for mode membership.  Two
nodes may have similar affinity vectors but be in complementary
learning states (one absorbing novelty while the other consolidates).
Phase coherence captures this dynamic alignment by comparing
\emph{surprisal velocities}:
\begin{equation}
  v_i = \bar{S}_i^{(t)} - \bar{S}_i^{(t-1)}
  \label{eq:phase_velocity}
\end{equation}
where $\bar{S}_i^{(t)}$ is node $i$'s mean surprisal at epoch $t$.
The phase coupling between two nodes is:
\begin{equation}
  C_{\text{phase}}(i, j)
    = \frac{v_i \cdot v_j}{\sqrt{|v_i|} \cdot \sqrt{|v_j|} + \epsilon}
  \label{eq:phase_coupling}
\end{equation}
with $\epsilon = 0.01$ to suppress spurious coupling when both nodes
are quiescent.  This is a normalized product: $C_{\text{phase}} = +1$
when both nodes are accelerating or both decelerating (in-phase),
$C_{\text{phase}} = -1$ when one is rising and the other falling
(anti-phase), and $C_{\text{phase}} \approx 0$ when at least one node
is stationary.

The geometric-mean normalization $\sqrt{|v_i| \cdot |v_j|}$ is used
instead of $\max(|v_i|, |v_j|)^2$ to avoid over-penalizing asymmetric
velocity pairs---a common case when one node has just begun learning
while a structurally aligned neighbor is in full swing.

The phase threshold for mode membership is $\theta_P = -0.3$, allowing
mild transient disagreement during mode formation.

\subsubsection{Eigenmode Detection}
\label{sec:eigenmode_detection}

Given the set of received digests $\{d_1, \ldots, d_N\}$, the
\texttt{FieldView} identifies eigenmodes by greedy clustering:

\begin{algorithm}[H]
\caption{Eigenmode Detection}
\label{alg:eigenmode}
\begin{algorithmic}[1]
\State $\mathcal{M} \leftarrow \emptyset$ \Comment{set of modes}
\For{each digest $d_i$}
  \State $m^* \leftarrow \arg\max_{m \in \mathcal{M}}
    C_{\text{aff}}(\text{center}(m), \mathbf{a}_i)$
    \textbf{s.t.}\ $C_{\text{aff}} \geq \theta_C$ and
    $C_{\text{phase}}(\text{seed}(m), d_i) \geq \theta_P$
  \If{$m^*$ exists}
    \State Add $d_i$ to $m^*$;
           update center: $\text{center}(m^*)_w \leftarrow
           \text{center}(m^*)_w \;\mathbin{|}\; \mathbf{a}_{i,w}$
    \State $E(m^*) \leftarrow E(m^*) + \bar{S}_i + |v_i|$
  \Else
    \State Create new mode $m = \{d_i\}$ with
           $\text{center}(m) = \mathbf{a}_i$,
           $E(m) = \bar{S}_i + |v_i|$
    \State $\mathcal{M} \leftarrow \mathcal{M} \cup \{m\}$
  \EndIf
\EndFor
\State \textbf{Dominant mode:}
       $m_{\text{dom}} = \arg\max_{m \in \mathcal{M}} E(m)$
\end{algorithmic}
\end{algorithm}

The energy of a mode $E(m) = \sum_{i \in m} (\bar{S}_i + |v_i|)$
aggregates both the absolute surprisal (how much the member nodes are
learning) and the magnitude of their learning velocity (how actively
they are changing).  The dominant mode---the one with the highest total
energy---is what the system ``attends to.''

Mode centers are computed incrementally via bitwise OR, which unions the
affinity bits of all members.  This makes mode boundaries inclusive:
a center grows to encompass its members rather than averaging them,
preventing the centroid-collapse problem of $k$-means in high-dimensional
binary spaces.

\subsubsection{Asymmetric Field Projection}

Once the dominant mode is identified, the \texttt{FieldView} projects
asymmetric pressure onto the owning node's MarkovTrie:

\begin{table}[h]
  \centering
  \caption{Field pressure multipliers.  The 4$\times$ decay differential
           and 5$\times$ learning differential produce selection by
           differential survival.}
  \label{tab:field_pressure}
  \begin{tabular}{lcc}
    \toprule
    & \textbf{Aligned (in dominant mode)}
    & \textbf{Misaligned (outside)} \\
    \midrule
    Decay multiplier   & $0.5\times$ (suppress) & $2.0\times$ (amplify) \\
    Learning multiplier & $1.5\times$ (amplify) & $0.3\times$ (suppress) \\
    \bottomrule
  \end{tabular}
\end{table}

These multipliers modulate the base field pressure, which is itself
derived from the weighted average of coupled neighbors' surprisal and
growth rates.  For a node with local surprisal $\bar{S}_{\text{local}}$
and field-averaged surprisal $\bar{S}_{\text{field}}$, the raw pressure
differentials are:
\begin{align}
  \Delta_{\text{decay}}  &= \bar{S}_{\text{field}} - \bar{S}_{\text{local}} \\
  \Delta_{\text{learn}}  &= \bar{S}_{\text{field}} - \bar{S}_{\text{local}} \\
  \Delta_{\text{prune}}  &= \dot{N}_{\text{field}} - \dot{N}_{\text{local}}
\end{align}
where the field averages are weighted by the product of affine coupling
and phase coupling:
$w_{ij} = C_{\text{aff}}(\mathbf{a}_i, \mathbf{a}_j) \cdot
(1 + C_{\text{phase}}(i, j))$.

The raw pressures are multiplied by the alignment-dependent factors from
\Cref{tab:field_pressure} and clamped to $[-3.0, +3.0]$ for decay and
learning (to prevent runaway feedback) and $[-5.0, +5.0]$ for pruning
(wider because pruning errors are less costly---a pruned node can regrow
from new data, while a corrupted decay factor distorts all future
predictions).

\subsubsection{Interpretation as Attention}

The Field mechanism implements attention without the machinery of
query-key-value projections.  The ``query'' is each node's current
adaptive state (broadcast as a digest).  The ``key matching'' is affine
coupling plus phase coherence.  The ``value aggregation'' is the
weighted field pressure from coupled neighbors.  The ``selection'' is
the asymmetric fork: amplify the attended mode, suppress the rest.

This produces three effects that are analogous to attention in
transformer architectures:

\begin{enumerate}[nosep]
  \item \textbf{Concentration}: aligned nodes retain knowledge longer
        (suppressed decay) and absorb related input faster (amplified
        learning), concentrating capacity on the dominant pattern.
  \item \textbf{Graceful transition}: because the suppression is
        multiplicative rather than binary, minority modes are dampened
        but not destroyed.  When the dominant pattern shifts, suppressed
        nodes still retain enough state to become the new dominant mode.
  \item \textbf{No catastrophic forgetting}: the 4$\times$ decay
        differential is strong enough to produce visible divergence
        within a few epochs but not so extreme that a temporarily
        misaligned node loses irrecoverable state (the decay multiplier
        is bounded at $2.0\times$, not $\infty$).
\end{enumerate}
